\documentclass[11pt]{article}
\usepackage[a4paper,margin=27mm]{geometry}
\usepackage{amsmath,amssymb,amsthm,booktabs,array}
\usepackage[T1]{fontenc}
\usepackage{lmodern}
\usepackage[colorlinks=true,linkcolor=blue,citecolor=blue,urlcolor=blue]{hyperref}
\newtheorem{theorem}{Theorem}
\newtheorem{corollary}[theorem]{Corollary}
\newcommand{\tr}{\operatorname{tr}}
\newcommand{\CM}{\mathcal M}
\newcommand{\Id}{\mathbb I}
\title{Refuting Peled et al.'s Section IV conjectures\\
after Theorems 1 and 3, Eqs.\ (10) and (12):\\
sharp fourth-order bounds without filter normal form}
\author{Haobo Ma\\\small ChronoAI Pte Ltd, Singapore
\and Wenlin Zhang\\\small National University of Singapore\\
\small ChronoAI Pte Ltd, Singapore}
\date{2 November 2026}
\begin{document}
\maketitle
\begin{abstract}
We refute the conjectures in Section IV immediately after Theorems 1 and 3,
Eqs.~(10) and (12), of Peled,
Te'eni, Carmi and Cohen, \emph{Scientific Reports} \textbf{11}, 2849
(2021), that filter normal form can be removed from their correlation-minor bounds.
An explicit separable state on $\mathbb C^2\otimes\mathbb C^4$ disproves
this extension at order $h=3$.
A single $3\times3$ minor gives a strict violation, even after adding
up to $1/200$ maximally mixed noise.
Every positive noise parameter in this interval gives a positive-definite
counterexample. The filter-normal-form theorems retain their stated scope.
A tetrahedral construction in dimensions $2\times D$, $5\le D\le8$,
also refutes both extensions at order four, for every Schatten parameter.
For every $D\ge4$, the sharp fourth-order separable bound on
$\mathbb C^2\otimes\mathbb C^D$ is $1/(192\sqrt3)$, attained by the
noiseless tetrahedral state. We also correct the tail constraint in
arXiv:2005.12079v2, Appendix D.5, Eq.~(D31).
\end{abstract}

\section{The question and its normalization}
Peled et al.~\cite{peled}, Section IV, conjecture immediately after
Theorem 3, Eq.~(12), that its hypothesis of filter normal form (FNF) can be removed;
Appendix D.5 of arXiv:2005.12079v2 gives supporting evidence.
We give a negative answer by explaining the rational witness encoded in
the \href{https://github.com/the-omega-institute/trureturing/blob/0f1f822c673c1903dc6b0a209a23dacebca2aa6c/D5/S3/Quantum/Entanglement/Witness/CorrelationMinorNormFNFRefutation.lean}{pinned Lean source},
and extend that witness to an interval of noise parameters by a written proof.
We also refute the parallel conjecture immediately after Theorem 1, Eq.~(10),
using four qubit projectors and orthogonal flags in a larger subsystem,
and prove the sharp fourth-order bound for every $D\ge4$.

For Hilbert--Schmidt orthonormal Hermitian bases $\{A_i\}$ and $\{B_j\}$,
use the correlation matrix
\[
 C_{ij}=\tr\bigl(\rho(A_i\otimes B_j)\bigr),\qquad
 \tr(A_iA_k)=\delta_{ik},\quad \tr(B_jB_l)=\delta_{jl}.
\]
Let $C_h(C)$ be the matrix of increasing-order $h\times h$ minors and
$\CM_{h,p}(\rho)=\|C_h(C)\|_p$, with Schatten norms, as in
\cite[Section III, Eqs.~(6)--(8)]{peled}. In particular $\|\cdot\|_\infty$
is the Euclidean operator norm, not the largest absolute entry.
Changing the local Hermitian bases multiplies $C$ by real orthogonal
matrices on the two sides: expand each new basis in the old one and use
the trace inner product. These transformations induce orthogonal maps
on the corresponding exterior powers, so the Schatten norms of $C_h(C)$
are unchanged. It is therefore legitimate to choose a diagonal matrix-unit
basis on the second subsystem instead of one containing a normalized identity.

Write $D=\max(d_A,d_B)$ and $d=\min(d_A,d_B)$.
Theorem 3 of \cite{peled} states, for separable FNF states and
$\sqrt{Dd}\le h\le d^2$, that
\begin{equation}\label{eq:bound}
 \CM_{h,\infty}(\rho)\le
 \frac1{\sqrt{Dd}}
 \left[\frac{D-1}{D(h-1)}\frac{d-1}{d(h-1)}\right]^{(h-1)/2}.
\end{equation}
FNF means that both reduced density matrices are maximally mixed.
For $d_A=2,d_B=4,h=3$, the parameter conditions read $\sqrt8<3\le4$
and the proposed upper bound is $3/(64\sqrt2)$.
The multipartite extension of Lenny et al.~\cite[Theorem 1]{lenny}
also retains FNF in its operator-norm bound.
Correlation-matrix separability criteria also appear in
de Vicente~\cite{devicente}; the issue here is the specific threshold
\eqref{eq:bound}, not the validity of correlation-based criteria in general.
Lupo, Aniello and Scardicchio~\cite{lupo} studied the elementary symmetric
polynomials and gave numerical separable bounds for $2\times2$ and
$2\times3$ systems. Their Tables 1--2 do not address $D\ge4$.
The covariance estimate of Zhang et al.~\cite[Theorem 1]{zhang}, also given
by Gittsovich et al.~\cite[Proposition IV.2]{covariance},
and the weighted trace-norm approach of Huang et al.~\cite[Theorem 1]{huang}
provide relevant general bounds. We give the short specialized arguments below.

\section{A rational family of separable counterexamples}
In the standard qubit basis define
\[
 P_0=\begin{pmatrix}1&0\\0&0\end{pmatrix},\quad
 P_1=\frac15\begin{pmatrix}1&2\\2&4\end{pmatrix},\quad
 P_2=\frac15\begin{pmatrix}1&-2\\-2&4\end{pmatrix}.
\]
Let $E_{jk}=|j\rangle\langle k|$ on $\mathbb C^4$, with indices
$0,1,2,3$, and set
\begin{equation}\label{eq:state}
 \rho_t=\frac{1-t}{3}\sum_{j=0}^2P_j\otimes E_{jj}
       +\frac{t}{8}\Id_8 .
\end{equation}

\begin{theorem}\label{thm:main}
For every $0\le t\le1/200$, the state $\rho_t$ is separable and satisfies
\[
 \CM_{3,\infty}(\rho_t)>\frac3{64\sqrt2}.
\]
For $t>0$ it is positive definite. None of these states is in FNF.
\end{theorem}

\begin{proof}
The three $P_j$ are the projectors onto the unit vectors
$(1,0)$, $(1,2)/\sqrt5$, and $(1,-2)/\sqrt5$.
Equation~\eqref{eq:state} is a convex combination of three pure product
states, with weights $(1-t)/3$, and the product state
$(\Id_2/2)\otimes(\Id_4/4)$, with weight $t$.
It is therefore separable, positive semidefinite, and has trace one.
Moreover $\rho_t\ge(t/8)\Id_8$, which gives positive definiteness when $t>0$.

On the qubit choose $A=(I,Z,X,Y)/\sqrt2$, with the usual Pauli matrices.
On $\mathbb C^4$ choose first $E_{00},E_{11},E_{22},E_{33}$,
then $(E_{jk}+E_{kj})/\sqrt2$ and
$(-iE_{jk}+iE_{kj})/\sqrt2$ for $j<k$.
These are complete Hermitian orthonormal bases: diagonal matrix units
are mutually orthogonal; each off-diagonal basis element has squared Hilbert--Schmidt norm
one; distinct supports are orthogonal; the symmetric and skew members
on the same support have trace product zero.

Put $a=(1-t)/3$, $r=(4-t)/12$, and $s=t/4$.
The correlation matrix has the form $C=2^{-1/2}[B\mid0]$, where
\begin{equation}\label{eq:matrix}
 B=\begin{pmatrix}
 r&r&r&s\\
 a&-3a/5&-3a/5&0\\
 0&4a/5&-4a/5&0\\
 0&0&0&0
 \end{pmatrix}.
\end{equation}
The expectations of $(I,Z,X,Y)$ in the three projectors are
\[
 (1,1,0,0),\qquad (1,-3/5,4/5,0),\qquad (1,-3/5,-4/5,0).
\]
The noise contributes $t/4$ to each diagonal
column of the identity row before multiplication by $2^{-1/2}$.
All off-diagonal Bob expectations vanish, since $\rho_t$ is block
diagonal in Bob's basis.

Take rows $I,Z,X$ and columns $E_{00},E_{11},E_{22}$.
The corresponding determinant is
\begin{equation}\label{eq:minor}
 m(t)=\frac1{(\sqrt2)^3}\frac{64ra^2}{25}
     =\frac1{2\sqrt2}\frac{16(4-t)(1-t)^2}{675}.
\end{equation}
To obtain $64ra^2/25$, factor $r,a,a$ from the three rows;
the remaining determinant is
$24/25+4/5+4/5=64/25$.
For any matrix $N$, $|N_{uv}|=|e_u^TN e_v|\le\|N\|_\infty$,
so $\CM_{3,\infty}(\rho_t)\ge m(t)$.
The inequality $m(t)>3/(64\sqrt2)$ is equivalent to
\begin{equation}\label{eq:scalar}
 512(4-t)(1-t)^2>2025.
\end{equation}
On $0\le t\le1/200$, both positive factors on the left are bounded
below by their values at $1/200$, and
\[
 512\frac{799}{200}\left(\frac{199}{200}\right)^2-2025
 =\frac{574}{15625}>0.
\]
This proves the strict violation throughout the interval.

Finally, Bob's reduced density matrix is
$\operatorname{diag}(r,r,r,s)$, with $r-s=(1-t)/3>0$.
It is not $\Id_4/4$, so $\rho_t$ is not in FNF.
\end{proof}

\section{The pinned witness and the scope of the conclusion}
The rational witness in the linked Lean source is $\rho_{1/1000}$:
its three pure-product weights are $333/1000$ and its noise term is
$\Id_8/8000$. At this parameter,
\[
 r=\frac{1333}{4000},\quad a=\frac{333}{1000},\quad
 s=\frac1{4000}.
\]
For completeness the four possibly nonzero entries of the selected
compound row, in column order $012,013,023,123$, are
\begin{equation}\label{eq:four}
 \frac1{(\sqrt2)^3}
 \left(\frac{64ra^2}{25},\ \frac{4sa^2}{5},\
       -\frac{4sa^2}{5},\ \frac{24sa^2}{25}\right).
\end{equation}
Every other compound row contains the zero $Y$ row, and any column
involving one of the last twelve columns of $C$ vanishes.
Thus $C_3(C)$ has one nonzero row and rank one; its operator norm
squared is exactly the sum of the four squares in \eqref{eq:four}.
This also explains why a row-energy estimate, as used in the formal
source, suffices.

\begin{corollary}
For $0\le t\le1/200$, all Schatten correlation-minor norms
$\CM_{3,p}(\rho_t)$, $1\le p\le\infty$, equal $\CM_{3,\infty}(\rho_t)$
and exceed $3/(64\sqrt2)$. The fourth-order correlation-minor norms vanish.
\end{corollary}
\begin{proof}
A rank-one matrix has exactly one nonzero singular value, so every
Schatten norm has that value. Equation~\eqref{eq:matrix} has rank at
most three, so all fourth-order minors vanish.
\end{proof}

This corollary does not refute the distinct $p=1$ threshold in
Theorem 1 of \cite{peled}. In dimensions $2\times4$, its third-order
right-hand side is $1/(16\sqrt2)+1/(48\sqrt6)$, rather than
$3/(64\sqrt2)$.

\section{The companion conjecture after Theorem 1}
Let $v_0,v_1,v_2,v_3$ be the four vectors
\[
 \frac{(1,1,1)}{\sqrt3},\quad
 \frac{(1,-1,-1)}{\sqrt3},\quad
 \frac{(-1,1,-1)}{\sqrt3},\quad
 \frac{(-1,-1,1)}{\sqrt3}.
\]
With $\sigma=(X,Y,Z)$, put $Q_j=(I+v_j\cdot\sigma)/2$.
Each $Q_j$ is a rank-one density projector: the Pauli anticommutation
relations give $(v_j\cdot\sigma)^2=I$, and $\tr(v_j\cdot\sigma)=0$.
For $D\ge4$ define
\[
 \omega_{D,t}=\frac{1-t}{4}\sum_{j=0}^3Q_j\otimes E_{jj}
             +\frac{t}{2D}\Id_{2D}.
\]
\begin{theorem}\label{thm:companion}
For every integer $5\le D\le8$ and $0\le t\le1/1000$, the separable
state $\omega_{D,t}$ violates the FNF-free extensions of both Theorem 1
and Theorem 3 of \cite{peled} at $h=4$.
It is positive definite for $t>0$. The violations hold for every
$1\le p\le\infty$ at the common fourth-order threshold.
\end{theorem}
\begin{proof}
The displayed expression is a convex mixture of product states, and
$\omega_{D,t}\ge t\Id_{2D}/(2D)$.
Its marginals are $\Id_2/2$ and
$\frac{1-t}{4}\sum_{j=0}^3E_{jj}+\frac{t}{D}\Id_D$;
the latter is nonuniform because $D>4$ and $t<1$.
Use $(I,X,Y,Z)/\sqrt2$ and a diagonal/off-diagonal basis on Bob's
system as before. The tetrahedral vectors satisfy
\[
 \sum_{j=0}^3 v_j=0,\qquad
 \sum_{j=0}^3v_jv_j^{\mathsf T}=\frac43\Id_3.
\]
Thus the correlation Gram matrix is diagonal with entries
\[
 q_0=\frac{(1-t)^2}{8}+\frac{2t-t^2}{2D},\qquad
 q_1=q_2=q_3=\frac{(1-t)^2}{24}.
\]
For example the identity row has first four entries
$[(1-t)/4+t/D]/\sqrt2$ and its remaining $D-4$ diagonal entries
are $t/(D\sqrt2)$, whose squared sum gives $q_0$.
The Pauli rows have entries $(1-t)v_{j,k}/(4\sqrt2)$,
which give the other three entries and the zero cross terms.
There are four positive singular values. Their product is the sole
nonzero singular value of $C_4(C)$, so for every $p$,
\[
 \CM_{4,p}(\omega_{D,t})^2
 =q_0\left(\frac{(1-t)^2}{24}\right)^3
 \ge\frac{(1-t)^8}{110592}.
\]
At $d=2,h=4$, the square of the right-hand side of \eqref{eq:bound} is
\[
 b_D^2=\frac{(D-1)^3}{11664D^4}.
\]
It decreases for real $D>4$, since its logarithmic derivative is
$(4-D)/(D(D-1))<0$. Bernoulli's inequality and $t\le1/1000$ give
\[
 \CM_{4,p}(\omega_{D,t})^2
 \ge\frac{1-8t}{110592}
 \ge\frac{124}{125\cdot110592}
 >\frac{64}{11664\cdot625}=b_5^2\ge b_D^2.
\]
For the strict step, the ratio of the left constant to $b_5^2$ is
$4185/4096>1$.
Theorem 3 applies to these parameters because $2D\le16=h^2$.
For Theorem 1, $D\le d^3=8$ and $h=4>1$. Its right-hand side
is $S_4(\alpha,\beta,\beta,\beta)=\alpha\beta^3$, where
$\alpha=1/\sqrt{2D}$ and $\beta=\sqrt{(D-1)/(18D)}$.
This is exactly $b_D$. Hence both conjectured extensions fail.
The fourth-order quantity is independent of $p$, as already shown.
\end{proof}

\section{Sharp replacement and neighboring consequences}

\begin{theorem}\label{thm:sharp}
For every separable state on $\mathbb C^2\otimes\mathbb C^D$, $D\ge2$,
and every $1\le p\le\infty$,
\[
 \CM_{4,p}(\rho)\le\frac1{192\sqrt3}.
\]
For each $D\ge4$, equality holds at $\omega_{D,0}$, so this is the
sharp bound in those dimensions without an FNF assumption.
\end{theorem}
\begin{proof}
Use the qubit basis $(I,X,Y,Z)/\sqrt2$ and any complete orthonormal
Hermitian Bob basis. Write a separable decomposition
$\rho=\sum_jw_j\alpha_j\otimes\beta_j$, with density matrices
$\alpha_j,\beta_j$ and nonnegative weights summing to one.
Let $a_j,b_j$ be their coefficient vectors, so $C=\sum_jw_ja_jb_j^T$.
The Bloch representation gives $a_j=(1,r_j)/\sqrt2$, $|r_j|\le1$,
and $\|b_j\|_2^2=\tr(\beta_j^2)\le1$.
Put $L=\operatorname{diag}(1,\sqrt3,\sqrt3,\sqrt3)$.
Since the nuclear norm of an outer product is the product of its
Euclidean norms, the triangle inequality gives
\[
 \|LC\|_1\le\sum_jw_j\|La_j\|_2\|b_j\|_2\le\sqrt2.
\]
If $\tau_1,\ldots,\tau_4$ are the four singular values of $LC$,
the arithmetic--geometric mean inequality implies
$\prod_i\tau_i\le(\sqrt2/4)^4=1/64$.
Also
\[
 \prod_i\tau_i=\sqrt{\det(LCC^TL^T)}
 =3\sqrt3\,\sqrt{\det(CC^T)}.
\]
The fourth compound has one row, so every Schatten norm equals
$\sqrt{\det(CC^T)}$, including the rank-deficient case where it vanishes.
This proves the bound. At $t=0$, the Gram entries of $\omega_{D,0}$ are
$1/8,1/24,1/24,1/24$; their product is $1/110592=(192\sqrt3)^{-2}$.
\end{proof}

\begin{corollary}\label{cor:design}
For $D\ge5$, no common-size pair of regular, coherent, degree-1 pure
designs in dimensions $2$ and $D$ exists, with the terminology of
\cite[Section IV]{peled}. Nevertheless $\omega_{D,0}$ attains the
unrestricted fourth-order separable maximum for every $p$.
\end{corollary}
\begin{proof}
Suppose a qubit design has $v$ pure projectors and common distinct-pair
overlap $\mu$. Coherence and trace give
$\sum_jP_j=(v/2)I_2$ and $\mu=(v-2)/(2(v-1))$.
The case $v=1$ cannot be coherent. For $v\ge2$, $0\le\mu<1$, so the
Hilbert--Schmidt Gram matrix $(1-\mu)I_v+\mu\mathbf1\mathbf1^T$
has rank $v$. The real space of Hermitian qubit matrices has dimension
four; hence $v\le4$. On Bob's side coherence requires
$\sum_jR_j=(v/D)I_D$, whose rank $D$ is at most $v$. Thus $D\le v\le4$.
The maximality assertion follows from Theorem~\ref{thm:sharp}.
\end{proof}
This determines the maximal-order qubit case adjacent to the
design-maximization hypothesis after Theorem 4: tetrahedral projectors
and orthogonal flags attain the optimum, although coherent paired designs
cannot do so for $D\ge5$. The broader hypothesis uses the term
``quantum designs'' without a precise characterization; it is not refuted
by the absence of coherent pairs. Other orders and dimensions are not determined here.

\begin{corollary}\label{cor:tail}
For any separable state on $\mathbb C^{d_A}\otimes\mathbb C^{d_B}$,
with $d=\min(d_A,d_B)$ and $D=\max(d_A,d_B)$, its correlation singular
values $s_1\ge\cdots\ge s_{d^2}$ satisfy
\[
 \sum_{i=2}^{d^2}s_i
 \le\sqrt{(1-\tr\rho_A^2)(1-\tr\rho_B^2)}
 \le\sqrt{\frac{d-1}{d}\frac{D-1}{D}}.
\]
Equation (D31) in \cite{peled}, arXiv v2, omits the square root on
its last expression and is false as written, even in FNF.
\end{corollary}
\begin{proof}
In a separable decomposition use coefficient vectors $a_j,b_j$ as above,
and their means $a,b$. Then
$K=C-ab^T=\sum_jw_j(a_j-a)(b_j-b)^T$.
The nuclear-norm triangle inequality and Cauchy--Schwarz give
\[
 \|K\|_1\le
 \sqrt{\Bigl(\sum_jw_j\|a_j\|^2-\|a\|^2\Bigr)
             \Bigl(\sum_jw_j\|b_j\|^2-\|b\|^2\Bigr)}
 \le\sqrt{(1-\tr\rho_A^2)(1-\tr\rho_B^2)}.
\]
The rank-one approximation inequality gives
$\sum_{i=2}^{d^2}s_i(C)\le\|C-ab^T\|_1$.
For completeness, if $P$ projects orthogonally to the span of $a$,
singular-value interlacing gives $s_{i+1}(C)\le s_i(PC)$, while
$PC=PK$ and $\|PK\|_1\le\|K\|_1$.
If $a=0$, take $P=I$ instead. Finally $\tr\rho_A^2\ge1/d_A$ and
$\tr\rho_B^2\ge1/d_B$.
For $\omega_{4,0}$, both marginals are maximally mixed, while its
tail sum is $3/\sqrt{24}=\sqrt{3/8}>3/8$, the bound printed in (D31).
The corrected inequality is saturated by this state.
\end{proof}
The covariance estimate in this proof is the separable-state inequality
of \cite[Theorem 1]{zhang} and \cite[Proposition IV.2]{covariance};
the proof here fixes its normalization.

\begin{corollary}\label{cor:equal}
For $d_A=d_B=d\ge2$ and $d\le h\le d^2$, the operator bound
\eqref{eq:bound} holds for every density matrix satisfying the CCNR
inequality $\sum_i s_i\le1$, without an FNF hypothesis. Thus in this
range the operator bound cannot detect entanglement beyond CCNR.
At $h=d^2$ this conclusion holds for every Schatten parameter.
\end{corollary}
\begin{proof}
Choose normalized identities as the first local observables. Then
$C_{00}=1/d$, so $x=s_1\ge1/d$, as in \cite[Lemma 2.1]{lipoon}.
AM--GM and CCNR give
\[
 \CM_{h,\infty}=\prod_{i=1}^h s_i
 \le x\left(\frac{1-x}{h-1}\right)^{h-1}
 \le\frac1d\left(\frac{d-1}{d(h-1)}\right)^{h-1}.
\]
For the second inequality, the function
$x(1-x)^{h-1}$ has derivative $(1-x)^{h-2}(1-hx)\le0$
on $1/d\le x\le1$, since $h\ge d$. The last expression is exactly
\eqref{eq:bound} when $D=d$. At $h=d^2$ the compound has a single
singular value, so all its Schatten norms coincide.
\end{proof}
The corresponding equal-dimensional $p=1$ CCNR comparison is already
proved by Li, Poon and Sze~\cite[Theorem 2.4]{lipoon}.

The neighboring source statements have the following coverage.
\begin{center}\small
\begin{tabular}{p{.34\linewidth}p{.57\linewidth}}
\toprule
Source item in \cite{peled} & Coverage and scope here\\
\midrule
Theorems 1 and 3, Eqs.~(10), (12), with FNF & Not covered as new claims; the original FNF theorems keep their scope.\\
Conjecture after Theorem 3 & Refuted-and-corrected: Theorem~\ref{thm:main}; sharp fourth-order replacement in Theorem~\ref{thm:sharp}.\\
Conjecture after Theorem 1 & Refuted-and-corrected: Theorem~\ref{thm:companion} for $5\le D\le8$, $h=4$.\\
Theorems 2 and 4, coherent designs & Prior proof cited: \cite[Appendices D.3, D.6]{peled}; their FNF hypotheses remain essential.\\
General-$p$ design hypothesis after Theorem 4 & Proved: sharp maximal-order qubit bound and coherent-pair obstruction (Theorem~\ref{thm:sharp}, Corollary~\ref{cor:design}); broader characterization not covered.\\
Appendix D.5, displayed Eq.~(D31) & Refuted-and-corrected: Corollary~\ref{cor:tail}; the covariance inequality has a prior proof in \cite{covariance}.\\
Section IV, equal-dimensional CCNR comparison & Proved for $p=\infty$, $d\le h\le d^2$, and every $p$ at $h=d^2$ (Corollary~\ref{cor:equal}); prior proof cited for $p=1$ \cite[Theorem 2.4]{lipoon}; other cases not covered.\\
Section V and Theorem 5, entropy and discord & Not covered.\\
\bottomrule
\end{tabular}
\end{center}

\paragraph{Verification scope and data.}
In the \href{https://github.com/the-omega-institute/trureturing/blob/0f1f822c673c1903dc6b0a209a23dacebca2aa6c/D5/S3/Quantum/Entanglement/Witness/CorrelationMinorNormFNFRefutation.lean}{linked Lean file},
the declaration
\begin{center}\small\ttfamily
D5.S3.Quantum.Entanglement.Witness.\\
CorrelationMinorNormFNFRefutation.result
\end{center}
proves $\neg\texttt{claim}$: the negation of the universally quantified
bound without FNF, using the fixed rational witness above.
The noise interval, the exact rank-one norm formula, the corollary,
Theorems~\ref{thm:companion} and \ref{thm:sharp}, and
Corollaries~\ref{cor:design}, \ref{cor:tail} and \ref{cor:equal}
are established by the written arguments in this note.
The accompanying \texttt{verify.py} reconstructs the fixed state and its
correlations with exact rational arithmetic, checks all third-order minors,
and checks the endpoint arithmetic in \eqref{eq:scalar}.
It also checks the tetrahedral sign identities and the fourth-order
inequalities for $D=5,6,7,8$ at $t=0,1/1000$.
It checks the sharp squared bound, the weighted-Gram equality at zero
noise for $4\le D\le12$, the dimension obstruction for coherent pairs,
and the squared tail comparison in Eq.~(D31).
The \LaTeX{} source and this finite verification program accompany the note.

\begin{thebibliography}{9}
\bibitem{peled}
B.~Y. Peled, A. Te'eni, A. Carmi and E. Cohen,
Correlation minor norms, entanglement detection and discord,
\emph{Scientific Reports} \textbf{11}, 2849 (2021).
\href{https://doi.org/10.1038/s41598-021-82303-3}{doi:10.1038/s41598-021-82303-3}.
Detailed appendices: \href{https://arxiv.org/abs/2005.12079v2}{arXiv:2005.12079v2}.
\bibitem{lenny}
R. Lenny, A. Te'eni, B.~Y. Peled, A. Carmi and E. Cohen,
Multipartite entanglement detection via correlation minor norm,
\emph{Quantum Information Processing} \textbf{22}, 292 (2023).
\href{https://doi.org/10.1007/s11128-023-04046-6}{doi:10.1007/s11128-023-04046-6}.
See also \href{https://arxiv.org/abs/2302.00650}{arXiv:2302.00650}, Theorem 1.
\bibitem{devicente}
J.~I. de Vicente,
Separability criteria based on the Bloch representation of density matrices,
\emph{Quantum Information and Computation} \textbf{7}, 624--638 (2007).
\href{https://arxiv.org/abs/quant-ph/0607195}{arXiv:quant-ph/0607195}.
\bibitem{lupo}
C. Lupo, P. Aniello and A. Scardicchio,
Bipartite quantum systems: on the realignment criterion and beyond,
\emph{Journal of Physics A: Mathematical and Theoretical} \textbf{41},
415301 (2008).
\href{https://doi.org/10.1088/1751-8113/41/41/415301}{doi:10.1088/1751-8113/41/41/415301}.
\bibitem{covariance}
O. Gittsovich, O. G\"uhne, P. Hyllus and J. Eisert,
Unifying several separability conditions using the covariance matrix criterion,
\emph{Physical Review A} \textbf{78}, 052319 (2008).
\href{https://doi.org/10.1103/PhysRevA.78.052319}{doi:10.1103/PhysRevA.78.052319}.
\href{https://arxiv.org/abs/0803.0757v3}{arXiv:0803.0757v3}.
\bibitem{zhang}
C.-J. Zhang, Y.-S. Zhang, S. Zhang and G.-C. Guo,
Entanglement detection beyond the computable cross-norm or realignment criterion,
\emph{Physical Review A} \textbf{77}, 060301(R) (2008).
\href{https://doi.org/10.1103/PhysRevA.77.060301}{doi:10.1103/PhysRevA.77.060301}.
\href{https://arxiv.org/abs/0709.3766v6}{arXiv:0709.3766v6}.
\bibitem{huang}
X. Huang, T. Zhang and N. Jing,
A unifying separability criterion based on extended correlation tensor (2024).
\href{https://arxiv.org/abs/2406.17230}{arXiv:2406.17230}, Theorem 1.
\bibitem{lipoon}
C.-K. Li, Y.-T. Poon and N.-S. Sze,
A note on the realignment criterion,
\emph{Journal of Physics A: Mathematical and Theoretical} \textbf{44},
315304 (2011).
\href{https://doi.org/10.1088/1751-8113/44/31/315304}{doi:10.1088/1751-8113/44/31/315304}.
\href{https://arxiv.org/abs/1103.2499v2}{arXiv:1103.2499v2}.
\end{thebibliography}
\noindent\small Copyright 2026 Haobo Ma and Wenlin Zhang.
Licensed under \href{https://creativecommons.org/licenses/by/4.0/}{CC BY 4.0}.
\end{document}
